% Part: history
% Chapter: biographies
% Section: rozsa-peter

\documentclass[../../../include/open-logic-section]{subfiles}

\begin{document}

\olfileid[hi]{his}{bio}{pet} 

\olsection{रोज़ा पीटर}

\olphoto{peter-rozsa}{रोज़ा पीटर}
 
रोज़ा पीटर का जन्म 17 फ़रवरी 1905 को हंगरी के बुडापेस्ट में रोज़ा पोलित्ज़र
के नाम से हुआ। वे पुनरावर्ती फलनों पर अपने कार्य के लिए सर्वाधिक प्रसिद्ध
हैं; पुनरावर्तन सिद्धांत के क्षेत्र की स्थापना में यह कार्य आवश्यक था।

पीटर का पालन-पोषण कठिन राजनीतिक समय में हुआ---उनकी किशोरावस्था में प्रथम
विश्व युद्ध चल रहा था---फिर भी वे बुडापेस्ट के संपन्न मारिया तेरेज़िया बालिका
विद्यालय जा सकीं, जहाँ से 1922 में स्नातक हुईं। फिर उन्होंने बुडापेस्ट के
पाज़मानी पीटर विश्वविद्यालय (जिसका नाम बाद में लोरांड एटवश विश्वविद्यालय
हुआ) में अध्ययन किया। पिता के आग्रह पर उन्होंने रसायन विज्ञान पढ़ना आरंभ
किया, किंतु बाद में गणित चुना और 1927 में स्नातक हुईं। यद्यपि वे माध्यमिक
विद्यालय में गणित पढ़ाने की योग्य थीं, महामंदी से विश्व अर्थव्यवस्था
प्रभावित होने के कारण उस समय आर्थिक स्थिति भयावह थी। इस दौरान पीटर ने गणित
की शिक्षिका और निजी अध्यापिका के रूप में छोटे-मोटे काम किए। अंततः वे गणित
में स्नातकोत्तर अध्ययन के लिए विश्वविद्यालय लौटीं। मूलतः उनकी योजना संख्या
सिद्धांत पर काम करने की थी, किंतु यह जानकर कि उनके परिणाम पहले ही सिद्ध किए
जा चुके हैं, उन्होंने लगभग गणित ही छोड़ दिया। उन्हें ग्योडेल के अपूर्णता
प्रमेयों पर काम करने के लिए प्रोत्साहित किया गया और उन्होंने अनजाने में उनके
कई परिणाम भिन्न तरीकों से सिद्ध कर दिए। इससे उनका आत्मविश्वास लौटा और डेविड
हिल्बर्ट के आधारभूत कार्यक्रम से प्रेरित पीटर ने पुनरावर्तन सिद्धांत पर अपने
पहले शोधपत्र लिखे। उन्होंने 1935 में पीएचडी प्राप्त की और 1937 में
\emph{Journal of Symbolic Logic} की संपादक बनीं।

पीटर के आरंभिक शोधपत्रों को व्यापक रूप से पुनरावर्ती फलन सिद्धांत के क्षेत्र
में संस्थापक योगदान माना जाता है। \cite{Peter1935a} में उन्होंने विभिन्न
प्रकार के पुनरावर्तन के बीच संबंध की जाँच की। \cite{Peter1935b} में उन्होंने
दिखाया कि एक निश्चित पुनरावर्ती रूप से परिभाषित फलन आदिम पुनरावर्ती नहीं है।
इससे विल्हेल्म एकरमान का पहले का परिणाम सरल हुआ। पीटर के सरलीकृत फलन को अब
प्रायः एकरमान फलन---और कभी-कभी अधिक उचित रूप से एकरमान--पीटर फलन---कहते हैं।
उन्होंने पुनरावर्ती फलन सिद्धांत पर पहली पुस्तक लिखी \citep{Peter1951}।

अपने कार्य के महत्त्व और प्रभाव के बावजूद पीटर को 1945 तक पूर्णकालिक
अध्यापन पद नहीं मिला। द्वितीय विश्व युद्ध में हंगरी के नाज़ी अधिग्रहण के
दौरान यहूदी-विरोधी कानूनों के कारण पीटर को पढ़ाने की अनुमति नहीं थी। 1944
में सरकार ने बुडापेस्ट में यहूदी बस्ती बनाई; उसे शेष नगर से काटकर सशस्त्र
पहरेदार तैनात किए गए। 1945 में बस्ती की मुक्ति तक पीटर को वहीं रहने के लिए
विवश किया गया। फिर उन्होंने बुडापेस्ट टीचर्स ट्रेनिंग कॉलेज और 1955 से एटवश
लोरांड विश्वविद्यालय में पढ़ाया। वे गणित की अकादमिक डॉक्टर बनने वाली पहली
हंगेरियाई महिला गणितज्ञ और हंगेरियन अकैडमी ऑफ साइंसेज़ के लिए निर्वाचित पहली
महिला थीं।

पीटर गणित की एक उत्साही शिक्षिका के रूप में जानी जाती थीं, जो केवल व्याख्यान
देने की अपेक्षा अपने विद्यार्थियों के साथ गणितीय समस्याओं की प्रकृति और
सुंदरता खोजना पसंद करती थीं। फलतः विद्यार्थी उन्हें स्नेह से ``आंट रोज़ा''
कहते थे। पीटर का निधन 1977 में~71 वर्ष की आयु में हुआ।

\begin{reading}
अधिक जीवनीपरक पठन के लिए \citep{Oconnor2014} और \citep{Andrasfai1986}
देखें। \citet{Tamassy1994} ने पीटर का एक संक्षिप्त साक्षात्कार लिया। गणित पर
मनोरंजक पठन के लिए पीटर की पुस्तक \emph{Playing With Infinity}
\citep{Peter2010} देखें।
\end{reading}

\end{document}
